\documentclass{DTETI_CP_C501}

% Penambahan Package yang akan Digunakan
\usepackage{graphicx}
\usepackage{setspace}
\usepackage{subfigure}
\usepackage{array}
\usepackage{indentfirst}
\usepackage{titlesec}
\usepackage{lipsum}
\usepackage[utf8]{inputenc}
\usepackage{multirow}
\usepackage{multicol}
\usepackage{caption}
\usepackage{longtable} 
\usepackage{algorithm}
\usepackage{algpseudocode}
% \usepackage[table,xcdraw]{xcolor}

\bibliographystyle{plain} 
\usepackage[numbers,sort&compress]{natbib}

% ------------------------------------------------------------ %
% Awal Dokumen
\begin{document}

% ------------------------------------------------------------ %
% Data Capstone
% Judul Capstone
\judul{Prototipe Sistem Pemantauan Kualitas Udara Kabin Kendaraan Berbasis Sensor CO dan CO\textsubscript{2} dengan Peringatan dan Aktuasi Otomatis}
\title{Prototype of an In-Vehicle Cabin Air Quality Monitoring System Based on CO and CO\textsubscript{2} Sensors with Automatic \mbox{Warning} and Actuation Mechanism}

% Jenis Dokumen
%% Jenis Dokumen : IMPLEMENTASI SOLUSI

% Kode Dokumen
%% Kode Dokumen : C-501

% Nomor Dokumen (Kode Tim/Kelompok Capstone)
\NoDok{B\_09}

% Nomor Revisi
\NoRev{00}

% Tanggal Penerbitan Dokumen
%% Otomatis terisi tanggal ketika file LaTeX ini di-compile

% Data Mahasiswa Capstone
%% Ketua kelompok : \MHS{<Nama Lengkap>}{<NIM>}{<Prodi>}{<Alamat Email>}
\MHS{Aurellya Ratna Dewanti}{23/517176/TK/56870}
	  {Teknologi Informasi}{aurellyaratnadewanti@mail.ugm.ac.id}

%% Anggota kelompok : \Anggota{<Nomor 1-3>}{<Nama Lengkap>}{<NIM>}{<Prodi>}{<Alamat Email>}
\Anggota{1}{Joecelyn Aurora Majesty}{23/514716/TK/56510}
	  {Teknologi Informasi}{joecelynauroramajesty@mail.ugm.ac.id}
\Anggota{2}{Maritza Vania Adelia}{23/517643/TK/56944}
	  {Teknologi Informasi}{maritzavaniaadelia@mail.ugm.ac.id}
\Anggota{3}{Muhammad Akbar Rabbani}{23/515815/TK/56739}
	  {Teknik Elektro}{muhammadakbarrabbani@mail.ugm.ac.id}

% Teks "PROGRAM STUDI" pada halaman judul (mahasiswa berasal dari dua program studi)
\ProdiJudul{Teknologi Informasi\\dan Teknik Elektro}

% Dosen Pembimbing
%% Format : {<Nama Lengkap>}{<NIP/NIU>}
\DPA{Dr. Ir. I Wayan Mustika, S.T., M.Eng.}{198109212014041001}

% Tempat Pelaksanaan
%% Format : {<Nama Laboratorium> \newline <Nama Departemen> \newline <Nama Fakultas>}
\Tempat{Departemen Teknik Elektro dan Teknologi Informasi \newline Fakultas Teknik}

% ------------------------------------------------------------ %
% Membuat Halaman Judul hingga Catatan Revisi Dokumen
\maketitle

% ------------------------------------------------------------ %
% Isi Laporan

% BAB 01 : Pengantar 
\chapter{\uppercase{Pengantar}}
\label{chap:Pengantar}
\section{Latar Belakang}
\label{sec:Latar_Belakang}

    % TODO (teman): isi subbab 1.1 Latar Belakang. Placeholder di bawah dihapus setelah diisi.
    \lipsum[1-3]

\section{Kendala-Kendala yang Perlu Dipertimbangkan}
\label{sec:Kendala_Kendala_yang_Perlu_Dipertimbangkan}

    Pengembangan AEROS dibatasi oleh kendala pada tiga aspek berikut.

    \begin{enumerate}
        \item \textbf{Aspek Teknis}\par
        Pada proyek ini, sensor gas, mikrokontroler, komunikasi nirkabel, basis data, dan aplikasi diintegrasikan agar bekerja di dalam kabin mobil. Kendala utamanya adalah keakuratan data sensor karena sensor gas semikonduktor juga merespons gas pereduksi lain selain CO dan mengalami \textit{drift} \cite{mics5524, mox_drift}. Kalibrasi dengan alat ukur referensi tidak dapat dilakukan sehingga konsentrasi CO hanya berupa estimasi. Sinyal sensor yang kecil menuntut pembacaan yang teliti dan penyaringan agar tidak memicu peringatan palsu. Komunikasi data \textit{real-time} bergantung pada jaringan Wi-Fi yang tidak selalu stabil. Getaran, suhu kabin, catu daya mobil, dan kelistrikan kendaraan untuk penggerak jendela juga harus diperhatikan.
        \item \textbf{Aspek Ekonomi}\par
        % TODO: "pengeluaran meningkat" bersumber dari pengguna; total belum dihitung ulang (hasil-dan-pengujian.md no. 8). Hitung dari nota.
        Pelaksanaan proyek ini harus memperhatikan biaya pengadaan komponen perangkat keras. Sistem dirancang dengan biaya yang efisien tanpa mengorbankan fungsi utamanya, yaitu pemantauan dan peringatan dini kualitas udara kabin. Banyak komponen berubah atau bertambah dari rencana awal selama implementasi sehingga pengeluaran meningkat. Batasan biaya ini memengaruhi pilihan komponen pada setiap subsistem. Harga seluruh komponen perlu dihitung ulang dan setiap penambahan komponen ditimbang terhadap manfaatnya bagi kinerja sistem.
        \item \textbf{Aspek Keselamatan dan Keamanan}\par
        % TODO: kalimat privasi belum punya sumber di aeros-sensor; konfirmasi dengan teman backend.
        Sistem digunakan di ruang yang dihuni manusia sehingga keselamatan pengemudi dan penumpang harus terjamin. Sistem berfungsi sebagai peringatan dini dan bukan alat ukur keselamatan bersertifikat. Tim tidak membuat gas CO tanpa protokol keselamatan sehingga kemampuan deteksi CO belum dapat diuji. Penggerak jendela memerlukan kehati-hatian karena fitur anti-jepit tidak dibangun dan kelistrikan mobil 12 V harus dirangkai dengan aman. Data yang dikirim harus terlindung dari pemalsuan dan penyadapan. Waktu pengukuran dan identitas perangkat juga harus dijaga privasinya.
    \end{enumerate}

\section{Mata Kuliah dan Kemampuan yang Mendukung}
\label{sec:Mata_Kuliah_dan_Kemampuan_yang_Mendukung}

Dalam merealisasikan sistem deteksi dan mitigasi gas beracun di dalam kabin kendaraan, proyek Capstone ini mengintegrasikan fondasi ilmiah dan keilmuan teknik yang telah dipelajari melalui mata kuliah berikut:

\begin{enumerate}[label=\alph*.]

    \item \textbf{Fisika Dasar}

    Mata kuliah ini mendasari analisis perilaku gas CO dan CO$_2$ sebagai fluida dalam ruang tertutup. Prinsip perpindahan kalor diaplikasikan untuk menghitung pengaruh perubahan suhu kabin yang ekstrem terhadap laju difusi gas dan akurasi sensor.

    \item \textbf{Statistika}

    \textit{Noise} elektronik dan variasi suhu sekitar merusak stabilitas sinyal. Ada juga efek \textit{cross-sensitivity} terhadap keberadaan gas interferensi. Konsep statistika menopang penentuan ambang batas peringatan yang tahan terhadap berbagai fluktuasi lingkungan tersebut.

    \item \textbf{Sistem Tertanam dan Internet of Things (IoT)}

    Prinsip mikrokontroler membentuk fondasi perangkat keras utama dari keseluruhan arsitektur pemantauan. Perancangan simpul fisik memanfaatkan konfigurasi pin cip ESP32 untuk mengintegrasikan pembacaan sensor menuju jaringan internet nirkabel.

    \item \textbf{Sistem Mikroprosesor}

    Arsitektur internal pemroses menuntut alokasi memori yang ketat dalam menangani banyak aliran data mentah secara bersamaan. Penentuan frekuensi pengambilan sampel gas diatur sedemikian rupa agar cip tetap responsif tanpa mengalami kelebihan beban.

    \item \textbf{Sistem Operasi pada Sistem Tertanam}

    Perangkat harus mengeksekusi banyak tugas kritis secara bersamaan tanpa saling menjegal antrean perintah. Kerangka kerja sistem operasi waktu nyata diterapkan pada mikrokontroler untuk mengatur prioritas eksekusi secara deterministik.

    \item \textbf{Pengukuran dan Instrumentasi}

    Prinsip instrumentasi mendasari strategi kalibrasi perangkat lunak untuk menekan kelemahan bawaan sensor gas berbiaya rendah. Perhitungan batas deviasi dan kompensasi suhu dirancang secara empiris agar peringatan bahaya tetap menyala pada waktu yang tepat.

    \item \textbf{Teknik Digital}

    Sinyal tegangan analog dari sensor diubah menjadi data biner melalui fitur konversi bawaan cip pemroses. Pemahaman tentang resolusi bit menentukan tingkat kepekaan deteksi kadar gas beracun pada rentang puluhan bagian per juta.

    \item \textbf{Jaringan dan Komunikasi Data}

    Arsitektur penyaluran pesan dari sensor menuju perangkat beroperasi mengikuti hierarki lapisan jaringan murni. Paradigma publikasi dan berlangganan diterapkan lewat standar protokol pertukaran pesan ringan untuk menjamin kelancaran lalu lintas data.

    \item \textbf{Sistem Komunikasi Nirkabel}

    Teori modulasi gelombang membenarkan pemilihan protokol jarak dekat hemat energi sebagai jembatan lalu lintas data dalam kabin. Kemampuan memancarkan paket informasi tanpa status koneksi permanen sangat menghemat daya operasi perangkat keras.

    \item \textbf{Teknik Pengolahan Isyarat Digital}

    Langkah penyaringan digital diterapkan pada baris kode untuk menghaluskan fluktuasi tajam akibat \textit{noise} elektromagnetik. Sinyal bersih hasil saringan ini mencegah perangkat mengambil keputusan aktuasi berdasarkan lonjakan angka sesaat.

    \item \textbf{Teknik Kendali}

    Logika pergerakan komponen mekanis menggunakan prinsip pengontrol hidup-mati dengan batas pita mati. Konsep histeresis membatasi motor jendela agar tidak bergerak membuka dan menutup secara liar saat kadar gas berada di ambang batas.

    \item \textbf{Pemrograman Mobile}

    Kemampuan ini digunakan untuk membangun aplikasi yang dapat menerima, menampilkan, dan mengolah data dari sensor (misalnya sensor CO, suhu, atau kelembapan) yang dikirim melalui jaringan nirkabel.

    \item \textbf{Pemrograman Dasar}

    Mata kuliah ini menjadi pondasi penting pemrograman baik untuk mikrokontroler maupun untuk \textit{back-end} dan \textit{front-end} website. Oleh karena itu, pemahaman dasar mengenai pemrograman sangat dibutuhkan untuk sistem monitoring.

    \item \textbf{Interaksi Manusia dan Komputer}

    Mata kuliah ini mendasari perancangan antarmuka (\textit{user interface}) aplikasi seluler yang intuitif agar penumpang dapat dengan cepat merespons peringatan darurat tanpa kebingungan.

\end{enumerate}


% BAB 02 : Dasar Teori Pendukung 
\chapter{\uppercase{Dasar Teori Pendukung}}
\label{chap:Dasar_Teori_Pendukung}
% TODO: subbab 2.3 sampai 2.9 masih placeholder (lorem ipsum). Isi sesuai kerangka-c501.md; 2.3 dan 2.8 dikerjakan bersama teman backend dan software.

\section{Kualitas Udara Kabin dan Batas Paparan}
\label{sec:Kualitas_Udara_Kabin}

    Kualitas udara dalam ruang (\textit{indoor air quality}, IAQ) dipengaruhi banyak zat sekaligus, misalnya gas dan partikel dari berbagai sumber \cite{ashrae_co2}. Kabin kendaraan adalah ruang kecil yang tertutup saat jendela dinaikkan. Dua gas yang dipantau pada proyek ini adalah karbon monoksida (CO) dan karbon dioksida (CO\textsubscript{2}).

    \subsection{Karbon Monoksida (CO)}

    CO dapat masuk ke kabin dari gas buang kendaraan lain. Studi di Auckland memakai tiga kendaraan dengan pengaturan ventilasi yang berbeda \cite{dirks2018}. Puncak CO sekitar 170~ppm terjadi pada pengaturan udara luar saat kendaraan mengikuti kendaraan tua yang berasap dan tidak terawat.

    Kadar CO ditulis dalam ppm (bagian per sejuta volume) atau dalam mg/m\textsuperscript{3}. NIOSH memakai faktor konversi 1~ppm sekitar 1,15~mg/m\textsuperscript{3} \cite{niosh_co}. WHO membatasi kadar CO dalam ruang hingga 35~mg/m\textsuperscript{3} (sekitar 30~ppm) untuk paparan satu jam \cite{who_iaq}.

    \subsection{Karbon Dioksida (CO\texorpdfstring{\textsubscript{2}}{2})}

    CO\textsubscript{2} dihasilkan oleh napas penumpang. Kadarnya naik saat kabin tertutup dan pertukaran udaranya kecil (Subbab~\ref{sec:Model_Kesetimbangan_Gas}). Pengukuran pada 51 mobil ber-AC di Hong Kong mencatat 96\% mobil melewati 1.000~ppm saat berkendara \cite{barnes2018}.

    % Results dan keterbatasan Satish dicek langsung ke artikel PMC3548274 (13 Nov 2026): 22 peserta (18 berusia 18--29 tahun, mayoritas mahasiswa),
    % 2,5 jam per kadar, skor sebagian besar skala keputusan lebih rendah pada 2.500 ppm. Penulis menyatakan temuan perlu direplikasi.
    Satish dkk. menguji 22 peserta, yang sebagian besar mahasiswa, selama 2,5 jam pada kadar 600, 1.000, dan 2.500~ppm. Skor pengambilan keputusan pada sebagian besar skala lebih rendah di 2.500~ppm \cite{satish2012}. Penulisnya menyatakan bahwa temuan ini sulit digeneralisasi dan perlu direplikasi \cite{satish2012}. Tinjauan ASHRAE menilai bukti pengaruh langsung CO\textsubscript{2} pada kadar umum dalam ruang tidak konsisten \cite{ashrae_co2}.

    \subsection{Relevansi CO dan CO\texorpdfstring{\textsubscript{2}}{2} terhadap Kualitas Udara Kabin}

    CO dan CO\textsubscript{2} menggambarkan dua hal yang berbeda. CO\textsubscript{2} dihasilkan penumpang, sehingga kadarnya menunjukkan kurangnya udara segar yang masuk dan bukan bersihnya udara. ASHRAE menyebut CO\textsubscript{2} sebagai indikator ventilasi, bukan indikator IAQ secara menyeluruh. Alasannya, polutan dari bahan bangunan atau dari luar tidak berkorelasi dengan CO\textsubscript{2} \cite{ashrae_co2}.

    Napas tidak menghasilkan CO, sehingga kadar CO tidak mengikuti jumlah penumpang. CO diukur sendiri karena gas ini memiliki batas paparan tersendiri (Tabel~\ref{tab:Ch02_Batas_Paparan}). Proyek ini hanya mengukur CO dan CO\textsubscript{2}, sehingga tidak dapat menilai IAQ secara luas. Partikel, gas lain, suhu, dan kelembapan tidak diukur secara terpisah. Istilah kualitas udara kabin pada laporan ini karena itu hanya berarti kadar CO dan CO\textsubscript{2}. CO\textsubscript{2} menjadi penanda ventilasi kabin dan CO menjadi penanda gas buang yang masuk.

    \subsection{Batas Paparan}

    Beberapa lembaga menetapkan batas paparan CO dan CO\textsubscript{2}. Tabel~\ref{tab:Ch02_Batas_Paparan} merangkum batas yang dipakai sebagai acuan pada proyek ini.

    \begin{longtable}{|L{6.0cm}|L{9.0cm}|}
        \caption{Batas Paparan CO dan CO\textsubscript{2}}
        \label{tab:Ch02_Batas_Paparan}
        \vspace{-0.75em}\\
        \hline
        \textbf{Acuan} & \textbf{Nilai} \\ \hline
        \endfirsthead
        \hline
        \textbf{Acuan} & \textbf{Nilai} \\ \hline
        \endhead
        \multicolumn{2}{|l|}{\textbf{Karbon monoksida (CO)}} \\ \hline
        NIOSH, REL \cite{niosh_co}      & 35~ppm (40~mg/m\textsuperscript{3}), TWA sampai 10 jam \\ \hline
        NIOSH, plafon (C) \cite{niosh_co} & 200~ppm (229~mg/m\textsuperscript{3}) \\ \hline
        NIOSH, IDLH \cite{niosh_co}     & 1.200~ppm \\ \hline
        OSHA, PEL \cite{osha_co}        & 50~ppm (55~mg/m\textsuperscript{3}), TWA 8 jam \\ \hline
        WHO, 15 menit \cite{who_iaq}    & 100~mg/m\textsuperscript{3} ($\approx$ 87~ppm) \\ \hline
        WHO, 1 jam \cite{who_iaq}       & 35~mg/m\textsuperscript{3} ($\approx$ 30~ppm) \\ \hline
        WHO, 8 jam \cite{who_iaq}       & 10~mg/m\textsuperscript{3} ($\approx$ 9~ppm) \\ \hline
        \multicolumn{2}{|l|}{\textbf{Karbon dioksida (CO\textsubscript{2})}} \\ \hline
        NIOSH, REL \cite{niosh_co2}     & 5.000~ppm, TWA \\ \hline
        OSHA, PEL \cite{niosh_co2, ashrae_co2} & 5.000~ppm, TWA \\ \hline
        ASHRAE \cite{ashrae_co2}        & Tidak menetapkan batas CO\textsubscript{2} dalam ruang. Angka 1.000~ppm bukan batas ASHRAE. \\ \hline
    \end{longtable}

    REL adalah \textit{recommended exposure limit}, yaitu batas yang direkomendasikan NIOSH. PEL adalah \textit{permissible exposure limit}, yaitu batas yang diizinkan OSHA. IDLH adalah kadar yang langsung membahayakan jiwa atau kesehatan. TWA (\textit{time-weighted average}) adalah rata-rata berbobot waktu. Nilai ppm pada baris WHO dihitung dari mg/m\textsuperscript{3} dengan faktor konversi di atas.

    Semua batas pada Tabel~\ref{tab:Ch02_Batas_Paparan} dibuat untuk tempat kerja atau ruang dalam gedung, bukan untuk kabin kendaraan. Studi di Hong Kong juga menyebut belum ada pedoman kadar polutan di kabin \cite{barnes2018}. Karena itu proyek ini tidak mengacu pada satu lembaga saja dan tidak mengklaim memenuhi standar tertentu. Setiap ambang peringatan dipilih dari batas yang sesuai dengan tujuannya, dan dasarnya dijelaskan pada Bab~4.

    % Peraturan Indonesia (Permenaker 5/2018, Permenkes 48/2016) sengaja tidak dimasukkan karena angka batasnya belum terverifikasi
    % (PDF resmi berupa gambar pindai). Tambahkan hanya bila tabel resminya sudah dibaca langsung dan memuat batas CO dan CO2.
    % Bibitem permenaker5 dan permenkes48 dihapus dari Referensi.tex; data lengkapnya ada di daftar-pustaka.md [25] dan [26].

\section{Model Kesetimbangan Massa Gas dalam Kabin}
\label{sec:Model_Kesetimbangan_Gas}

    % Sumber: C-251 Bab 4 (PDF "[v1] C251 B-09.pdf"). Nomor persamaan di PDF: 4.1 (neraca), 4.2 (N), 4.2a (ringkas), 4.3 (sumber berhenti),
    % 4.4 (waktu), 4.5 (kesetimbangan, sumber berkelanjutan), 4.6 (P untuk CO2), 4.7 (kesetimbangan CO2).
    % Nomor di kerangka-c501.md dan di docstring pemodelan_bab4_simulasi.py (4.2, 4.3, 4.4a) berbeda dari PDF.
    % Sengaja tidak dimasukkan ke Bab 2 (pindahkan ke Bab 4, bagian implementasi): C-251 4.1.7 (adaptasi untuk estimasi CO2, hanya bentuk P yang dipakai di sini)
    % dan 4.1.8 (adaptasi untuk CO dan implikasinya pada logika status WASPADA).
    % TODO: bagian CO bersifat teori. Validasi CO dengan data uji tidak mungkin tanpa CO acuan (model-vs-uji.md bagian 5); sesuaikan Bab 5 bila prototipe CO berubah.

    Model kesetimbangan massa menggambarkan perubahan kadar gas di dalam kabin terhadap waktu. Perubahan itu dipengaruhi oleh sumber gas di dalam kabin dan pertukaran udara dengan luar kabin. Model ini berlaku untuk CO maupun CO\textsubscript{2}.

    \subsection{Model Satu Zona Tercampur Sempurna}

    Model ini memperlakukan kabin sebagai satu volume kontrol dengan udara tercampur merata (\textit{single-zone well-mixed}). Satu nilai kadar gas $C$ lalu mewakili seluruh kabin. Asumsi yang sama dipakai pada metode peluruhan untuk mengukur laju pertukaran udara \cite{ashrae_co2}. Model diturunkan dari hukum kekekalan massa, yaitu laju perubahan massa gas di dalam kabin sama dengan laju massa yang masuk dikurangi laju massa yang keluar.

    Gas masuk lewat dua jalur: sumber di dalam kabin dengan debit $Q_{\text{sumber}}$ dan kadar $C_{\text{sumber}}$, serta udara luar dengan debit $Q_{\text{inf}}$ dan kadar $C_{\text{luar}}$. Gas keluar bersama udara kabin. Volume kabin $V$ dianggap tetap dan $Q_{\text{sumber}}$ dianggap jauh lebih kecil daripada $Q_{\text{inf}}$, sehingga debit udara keluar sama dengan $Q_{\text{inf}}$. Karena $V$ tetap, laju perubahan massa gas sebanding dengan $V\,dC/dt$. Neraca massanya adalah
    \begin{equation}
        \label{eqn:Ch02_Neraca}
        V\,\frac{dC}{dt} = Q_{\text{sumber}}\,C_{\text{sumber}} + Q_{\text{inf}}\,C_{\text{luar}} - Q_{\text{inf}}\,C
    \end{equation}

    \subsection{Laju Pertukaran Udara dan Suku Sumber}

    Laju pertukaran udara (\textit{air changes per hour}, ACH) didefinisikan sebagai
    \begin{equation}
        \label{eqn:Ch02_N}
        N = \frac{Q_{\text{inf}}}{V}
    \end{equation}
    dalam satuan per jam. Besaran ini menyatakan berapa kali seluruh udara kabin diganti dalam satu jam. Dengan $N$ dan laju penambahan oleh sumber $P = Q_{\text{sumber}}\,C_{\text{sumber}}/V$ (ppm per jam), Persamaan~\ref{eqn:Ch02_Neraca} menjadi
    \begin{equation}
        \label{eqn:Ch02_Model_Gas}
        \frac{dC}{dt} = P - N\left(C - C_{\text{luar}}\right)
    \end{equation}
    Suku $P$ menaikkan kadar dengan laju tetap, sedangkan suku $N\left(C - C_{\text{luar}}\right)$ menggerakkan kadar kabin menuju kadar udara luar.

    Bentuk $P$ berbeda untuk tiap gas. Sumber CO\textsubscript{2} adalah napas penumpang:
    \begin{equation}
        \label{eqn:Ch02_P_CO2}
        P = \frac{n\,G\,10^{6}}{V}
    \end{equation}
    dengan $n$ jumlah penumpang dan $G$ volume CO\textsubscript{2} yang dihembuskan satu orang per jam (L/jam). Faktor $10^{6}$ mengubah perbandingan volume menjadi ppm. Sumber CO berupa sumber internal atau kebocoran dengan bentuk $P$ sebelumnya. CO dari luar kabin tidak masuk lewat $P$ melainkan lewat $C_{\text{luar}}$.

    \subsection{Penyelesaian untuk Sumber Berhenti}

    Bila sumber berhenti ($P = 0$), Persamaan~\ref{eqn:Ch02_Model_Gas} memiliki penyelesaian
    \begin{equation}
        \label{eqn:Ch02_Berhenti}
        C(t) = C_{\text{luar}} + \left(C_{0} - C_{\text{luar}}\right)e^{-Nt}
    \end{equation}
    dengan $C_{0}$ kadar awal dan $t$ waktu (jam). Kadar bergerak menuju kadar luar dari kedua arah. Kadar turun bila $C_{\text{luar}} < C_{0}$ dan naik bila $C_{\text{luar}} > C_{0}$. Jadi pembukaan jendela mengencerkan gas bila udara luar lebih bersih, tetapi justru memasukkan gas bila udara luar lebih tercemar.

    Waktu yang dibutuhkan kadar untuk mencapai target $C_{\text{target}}$ diperoleh dengan membalik Persamaan~\ref{eqn:Ch02_Berhenti}:
    \begin{equation}
        \label{eqn:Ch02_Waktu_Target}
        t = \frac{1}{N}\ln\frac{C_{0} - C_{\text{luar}}}{C_{\text{target}} - C_{\text{luar}}}
    \end{equation}
    Persamaan ini hanya berlaku bila $C_{\text{target}}$ berada di antara $C_{0}$ dan $C_{\text{luar}}$. Bila $C(t)$ terukur pada waktu $t$, persamaan yang sama dapat disusun ulang untuk memperkirakan laju pertukaran udara dengan metode peluruhan \cite{ashrae_co2}:
    \begin{equation}
        \label{eqn:Ch02_Estimasi_N}
        N = \frac{1}{t}\ln\frac{C_{0} - C_{\text{luar}}}{C(t) - C_{\text{luar}}}
    \end{equation}

    \subsection{Penyelesaian untuk Sumber Berkelanjutan}

    Bila sumber terus ada ($P > 0$), kadar menuju kadar kesetimbangan $C_{\infty}$ yang diperoleh dari $dC/dt = 0$:
    \begin{equation}
        \label{eqn:Ch02_Kesetimbangan}
        C_{\infty} = C_{\text{luar}} + \frac{P}{N}
    \end{equation}
    Kadar tidak pernah turun sampai $C_{\text{luar}}$. Nilai $C_{\infty}$ makin tinggi untuk sumber yang makin kuat atau $N$ yang makin kecil. Jendela terbuka memperbesar $N$ sehingga $C_{\infty}$ mendekati $C_{\text{luar}}$. Pada CO\textsubscript{2}, $C_{\infty}$ ditentukan oleh jumlah penumpang, volume kabin, dan $N$.

    Penyelesaian lengkapnya berbentuk
    \begin{equation}
        \label{eqn:Ch02_Penyelesaian_Umum}
        C(t) = C_{\infty} + \left(C_{0} - C_{\infty}\right)e^{-Nt}
    \end{equation}
    Waktu mencapai target memakai Persamaan~\ref{eqn:Ch02_Waktu_Target} dengan $C_{\text{luar}}$ diganti $C_{\infty}$. Persamaan itu berlaku bila $C_{\text{target}}$ berada di antara $C_{0}$ dan $C_{\infty}$.

    Model ini dipakai pada Bab~4 untuk memperkirakan lama kadar CO\textsubscript{2} dan CO mencapai ambang serta lama pemulihannya. Perhitungannya memakai skrip pemodelan tim \cite{skrip_pemodelan}, dan hasilnya dibandingkan dengan data uji pada Bab~5.

\section{Internet of Things}
\label{sec:Internet_of_Things}

    % TODO: placeholder. 2.3.1 MQTT (publish/subscribe, broker, QoS, TLS, HMAC-SHA256) dan 2.3.2 HTTP dan REST. Kerjakan bersama teman backend.
    \lipsum[1]

\section{Development Board: ESP32}
\label{sec:Development_Board_ESP32}

    % TODO: placeholder. mikrokontroler, Wi-Fi, ADC 12-bit, UART, I2C, GPIO; versi DevKitC V4; gambar papan dengan sumber.
    \lipsum[1]

\section{Sensor}
\label{sec:Sensor}

    % Sumber: parameter-dan-sumber.md (bagian A dan B), batasan-sistem.md, datasheet MiCS-5524 rev 6 [1], manual MH-Z19C v1.2 [2], Adafruit [4], model koreksi MOX [6].
    % Tidak dikutip: Liu dan Tang [7] (volume dan halaman belum dicek, belum dibaca langsung).
    % TODO: (1) gambar masih placeholder; ambil foto modul sendiri atau gambar dari sumber dengan izin, lalu ganti kotak di bawah.
    % (2) MQ-7 dan MH-Z19B sengaja tidak ditulis: pengembangan memakai MiCS-5524 dan MH-Z19C.
    % (3) Cek stiker modul MH-Z19C untuk memastikan varian rentang 5.000 ppm (peringatan-dan-ambang.md A6 no. 5).

    \subsection{Sensor Gas Semikonduktor (MOX)}
    \label{sec:Sensor_MOX}

    Sensor gas semikonduktor oksida logam (\textit{metal oxide semiconductor}, MOX) memiliki lapisan yang hambatannya ($R_s$) turun saat terkena gas pereduksi seperti CO \cite{mics5524}. Elemen pemanas di dalam sensor menjaga lapisan itu pada suhu kerja \cite{mics5524}.

    Hambatan $R_s$ dibaca lewat pembagi tegangan dengan resistor beban $R_L$. Bila $V_{CC}$ adalah tegangan catu dan $V_{out}$ tegangan pada $R_L$, maka
    \begin{equation}
        \label{eqn:Ch02_Rs}
        R_s = R_L\,\frac{V_{CC} - V_{out}}{V_{out}}
    \end{equation}
    Datasheet menyatakan kadar gas lewat rasio $R_s/R_0$, dengan $R_0$ hambatan sensor di udara bersih \cite{mics5524}. Kadar $c$ (ppm) dihitung dari rasio itu dengan hubungan pangkat
    \begin{equation}
        \label{eqn:Ch02_Pangkat}
        c = A\left(\frac{R_s}{R_0}\right)^{B}
    \end{equation}
    dengan $A$ dan $B$ konstanta dari kurva pada datasheet ($B$ negatif).

    Sensor MOX memiliki tiga keterbatasan. Pertama, sensor tidak selektif karena gas pereduksi lain seperti etanol, hidrogen, dan amonia ikut mengubah $R_s$ \cite{mics5524, adafruit_mics}. Kedua, $R_s$ bergeser terhadap waktu (\textit{drift}) serta dipengaruhi suhu dan kelembapan \cite{mox_drift}. Ketiga, kurva datasheet hanya berlaku pada 25~$^{\circ}$C dan 50\% kelembapan relatif \cite{mics5524}.

    MiCS-5524 adalah sensor MOX untuk CO, etanol, hidrogen, amonia, metana, propana, dan isobutana dengan rentang CO 1--1.000~ppm \cite{mics5524}. Pemanasnya bertegangan 2,4~V (76~mW) dengan catu 4,9--5,1~V, dan $R_0$ berkisar 100--1.500~k$\Omega$ antarunit \cite{mics5524}. Datasheet tidak memuat waktu respons. Modul yang dipakai berkaki tiga, yaitu 5~V, GND, dan keluaran analog (Gambar~\ref{fig:Ch02_MiCS5524}).

    \begin{figure}[h!]
        \centering
        \fbox{\parbox[c][3cm][c]{6cm}{\centering\textit{Gambar menyusul}}}
        \caption{Modul sensor MiCS-5524}
        \label{fig:Ch02_MiCS5524}
    \end{figure}

    \subsection{Sensor CO\texorpdfstring{\textsubscript{2}}{2} Inframerah Nondispersif (NDIR)}
    \label{sec:Sensor_NDIR}

    Sensor inframerah nondispersif (\textit{non-dispersive infrared}, NDIR) mengukur CO\textsubscript{2} dari penyerapan cahaya inframerah. Makin tinggi kadar CO\textsubscript{2}, makin besar cahaya yang diserap \cite{mhz19c}.

    MH-Z19C adalah modul NDIR yang mengirim kadar CO\textsubscript{2} dalam ppm lewat UART 9600 baud pada level 3,3~V \cite{mhz19c}. Varian 5.000~ppm memiliki rentang 400--5.000~ppm dengan akurasi $\pm$(50~ppm + 5\% pembacaan). Waktu respons T90, yaitu waktu mencapai 90\% perubahan kadar, kurang dari 120 detik dan pemanasan awalnya 1 menit \cite{mhz19c}. Kalibrasi otomatis (\textit{automatic baseline correction}, ABC) berjalan tiap 24 jam dengan titik nol 400~ppm. Fitur ini kurang cocok untuk ruang yang jarang terkena udara luar \cite{mhz19c} (Gambar~\ref{fig:Ch02_MHZ19C}).

    \begin{figure}[h!]
        \centering
        \fbox{\parbox[c][3cm][c]{6cm}{\centering\textit{Gambar menyusul}}}
        \caption{Modul sensor MH-Z19C}
        \label{fig:Ch02_MHZ19C}
    \end{figure}

\section{Komponen Antarmuka dan Aktuator}
\label{sec:Komponen_Antarmuka_Aktuator}

    % TODO: placeholder. 2.6.1 ADS1115; 2.6.2 LCD I2C 16x2 dan buzzer aktif; 2.6.3 relay dan penggerak jendela. Tiap komponen satu gambar.
    \lipsum[1]

\section{Sumber Energi DC}
\label{sec:Sumber_Energi_DC}

    % TODO: placeholder. catu mobil 12 V, buck converter ke 5 V, adaptor USB; pentingnya catu 5 V stabil untuk sensor MOX.
    \lipsum[1]

\section{Sistem Informasi}
\label{sec:Sistem_Informasi}

    % TODO: placeholder. Supabase (PostgreSQL, REST), worker MQTT, aplikasi Flutter. Ditulis teman backend dan software.
    \lipsum[1]

\section{Bahasa Pemrograman}
\label{sec:Bahasa_Pemrograman}

    % TODO: placeholder. C/C++ (Arduino framework dan PlatformIO), JavaScript (worker Node.js), Dart (Flutter; cek dengan teman software).
    \lipsum[1]


% BAB 03 : Analisis Studi Pustaka Kunci dan Pemilihan Metode
\chapter{\uppercase{Analisis Studi Pustaka Kunci dan Pemilihan Metode}}
\label{chap:Analisis_Studi_Pustaka_Kunci_Pemilihan_Metode}
Salah satu kemampuan yang harus dimiliki oleh seorang Sarjana Teknik adalah kemampuan menganalisis sumber-sumber pustaka. Oleh karena itu, di dalam dokuman C-251 ini, mahasiswa perlu melakukan studi pustaka dan menuliskannya sampai kepada level kedalaman dan detail yang memadai. Sebelum melakukan analisis pustaka, uraikan lagi apa saja permasalahan yang ingin Anda selesaikan di dalam \textit{capstone} ini. Dari permasalahan tersebut, carilah literatur yang sudah pernah membahas dan menyelesaikan masalah yang sama, dan pilihkan pustaka kunci yang akan Anda adopsi untuk membantu menyelesaikan Anda. 

Untuk melakukan analisis studi, tentukan minimal tiga pustaka yang paling berkaitan dengan solusi yang diusulkan. Buatlah analisis yang mendalam dan jika perlu disimulasikan metodenya untuk menyelesaikan permasalahan Anda. Uraikan secara ringkas sumber-sumber yang membahas berbagai kerangka teori dan metode-metode yang mungkin untuk digunakan secara sistematis. Gunakan \textit{heading} untuk menguraikan metode yang ada.

\section{Metode 1}
\label{sec:Metode_01}

    \lipsum[1-3]

\section{Metode 2}
\label{sec:Metode_02}

    \lipsum[4-6]
    
    \subsection{Sub-Metode 2}
    \label{subsec:Sub-Metode_02}
    
    \lipsum[7-9]

\section{Metode 3}
\label{sec:Metode_03}

    \lipsum[10-12]
    
    \subsection{Sub-Metode 3A}
    \label{subsec:Sub-Metode_03A}
    
    \lipsum[13-15]
    
    \subsection{Sub-Metode 3B}
    \label{subsec:Sub-Metode_03B}
    
    \lipsum[15-18]


\section{Pemilihan Metode}
\label{sec:Pemilihan_Metode}
    
    Setelah dideskripsikan kekurangan dan kelebihan metode, tim akan memilih metode yang paling tepat disertai pengembangan metode tersebut. Jika dari semua metode yang ada tidak menghasilkan solusi yang konkret, tim bisa membuat metode sendiri. Dalam hal ini, proyek \textit{capstone} ini akan menjalankan mode penelitian sehingga tidak diwajibkan untuk menghasilkan luaran tertentu. Walaupun begitu adanya luaran seperti tertulis pada lembar luaran sangat disarankan. 
    
    \textbf{\uline{Apabila metode yang diimplementasikan berbeda dengan apa yang Anda usulkan dalam dokumen C-251}}, maka berikan penjelasan detail mengenai perubahan metode tersebut. Sampaikanlah pula argumen/justifikasi mengapa perubahan metode tersebut perlu Anda lakukan dalam proses implementasi.

% BAB 04 : Detail Implementasi
\chapter{\uppercase{Detail Implementasi}}
\label{chap:Detail_Implementasi}

\section{Luaran \textit{Capstone Project} beserta Spesifikasinya}
\label{sec:Luaran_Spesifikasi}
    
    Pada bagian ini, mahasiswa menjelaskan luaran (\textit{output}) apa yang dijanjikan di akhir periode capstone. Bagian ini harus disinkronkan dengan lembar luaran pada awal dokumen ini. Ada berbagai jenis luaran yang harus dihasilkan minimal satu oleh mahasiswa seperti disebutkan pada lembar luaran.

    Pertama-tama, mahasiswa menjelaskan luaran (\textit{output}) apa yang dijanjikan di akhir periode \textit{capstone}. Bagian ini harus disinkronkan dengan lembar luaran pada awal dokumen ini. Ada berbagai jenis luaran yang harus dihasilkan minimal satu oleh mahasiswa seperti disebutkan pada lembar luaran seperti yang dapat dituliskan pada Tabel \ref{tab:Ch06_Contoh_Luaran}.
    
    \begin{longtable}{|L{7.5cm}|L{7.5cm}|}
        \caption{Contoh Luaran} 
        \label{tab:Ch06_Contoh_Luaran}
        \vspace{-0.75em}\\        
        \hline
        \multicolumn{1}{|c|}{\textbf{Jenis Luaran}}                                                                     & \multicolumn{1}{c|}{\textbf{Contoh}}                                                                                \\ \hline
        
        \textit{Hardware} Digital (fisik)                                                                                        & Alat pendeteksi asap berbasis mikroprosesor                                                                         \\ \hline
        
        \textit{Hardware} Analog (fisik)                                                                                         & Alat penghasil frekuensi 1 THz menggunakan rangkaian RLC                                                            \\ \hline
        
        \textit{Firmware}/SW di Mikroprosesor/\textit{Development Board}                                                                  & Alat untuk memonitor asap berbasis komputer                                                                         \\ \hline
        
        \textit{Software} di PC                                                                                                  & Software untuk mendeteksi asap berbasis Python                                                                      \\ \hline
        
        Sistem Informasi                                                                                                & Sistem informasi \textit{capstone} mahasiswa, Sistem informasi manajemen kebakaran                                           \\ \hline
        Simulasi Lengkap                                                                                                & Simulasi black-out dengan MATLAB, C++ atau Python (bukan dengan \textit{software} jadi misalnya ETAP atau Digisilent)        \\ \hline
        
        Prototipe/miniatur \textit{hardware}/\textit{software}/sistem                                                                     & Prototipe lampu lalu lintas cerdas perempatan MM-UGM                                                                \\ \hline
        
        Teorema/Teori Baru                                                                                              & Teori kestabilan sistem tenaga terbarukan dengan beban negatif resistif                                             \\ \hline
        
        Kebaruan/\textit{novelty} yang lain/perbaikan metode (apabila berupa \textit{capstone} penelitian)                                & Perbaikan faktor daya dengan metode Alpha, Pencarian jarak terpendek pada sensor asap                               \\ \hline
        
        Dokumen (Kebijakan, SOP, Lingkungan, Ekonomi, dll)                                                              & SOP teknis perencanaan sistem proteksi, Dokumentasi \textit{black-out} dan analisisnya, Dokumen usulan investasi energi, dll \\ \hline
        
        \begin{tabular}[C{7.5cm}]{@{}L{7.5cm}@{}}Lain-lain (sebutkan)\\ Contoh : Simulasi parsial (dengan \textit{software} jadi)\end{tabular} & Contoh : Simulasi trafo dengan Fluent                                                                            \\ \hline
    \end{longtable}

    Kemudian Anda jelaskan pula target spesifikasi luaran dari \textit{Capstone Project}ini. Target performa yang akan dicapai berdasarkan fungsionalitas sebagai syarat keberhasilan kinerja produk harus dicantumkan (misal : modul kontroller dapat mengurangi \textit{overshoot} hingga dibawah 1\% atau akurasi produk di bawah 0.01\%). Yang perlu ditekankan disini mahasiswa harus menggunakan standar-standar keteknikan (besaran dan satuan, \textit{policy}/aturan, dan lain-lain) yang disepakati oleh organisasi keteknikan misalnya IEEE, PUIL, ACM dan lain-lain. Contoh spesifikasi luaran yang baik adalah seperti pada Tabel \ref{tab:Ch06_Contoh_Spesifikasi_Luaran}.
    
    \begin{longtable}{|c|c|c|c|c|}
        \caption{Contoh Spesifikasi Luaran} 
        \label{tab:Ch06_Contoh_Spesifikasi_Luaran}
        \vspace{-0.75em}\\
        \hline
        \textbf{No} & \textbf{Spesifikasi}  & \textbf{Satuan} & \textbf{Standar} & \textbf{Keterangan} \\ \hline
        1           & Tegangan Masukan      & Volt (V)        & 95 sampai 220 V  & Lihat Penjelasan A  \\ \hline
        2           & Tegangan Keluaran     & Volt (V)        & 20 V $\pm$ 0,2\% & Lihat Penjelasan B  \\ \hline
        3           & Interferensi Magnetis & Watt (W)        & Maksimal 0,1 W   & Lihat Penjelasan C  \\ \hline
        4           & SNR                   & Decibel (dB)    & Minimum 80 dB    & Lihat Penjelasan D  \\ \hline
        5           & dan seterusnya        &                 &                  &                     \\ \hline
    \end{longtable}

\section{Batasan Masalah}
\label{sec:Batasan_Masalah}
    
    Pada bagian ini peserta \textit{capstone} menguraikan lebih lanjut spesifikasi dari rancangan produk di atas. Spesifikasi yang dimaksud di sini adalah spesifikasi yang bersifat fungsional. Tim juga harus mendefinisikan batasan-batasan permasalahan secara lebih terperinci dan konkret. Batasan-batasan ini diperlukan untuk membuat sistem yang dirancang bisa diimplementasikan secara masuk akal dalam jangka waktu, kesediaan sarana dan prasarana, jumlah sumber daya manusia (SDM), biaya yang terbatas, dan lain-lain. 

    Pada bidang TI, peserta \textit{capstone} dapat menjelaskan fasilitas atau fitur apa saja yang disediakan oleh sistem yang dirancang secara lebih detail. Sebagai contoh, sistem informasi keuangan dapat dispesifikasikan untuk memiliki kemampuan untuk melakukan perhitungan biaya atau gaji (\textit{payroll}) berserta laporan (\textit{reporting}). Pada kasus ini, peserta wajib menyebutkan secara detail jenis pelaporan yang ada. Kemudian, perlu juga dijelaskan bagaimana perhitungan tersebut dilakukan, serta \textit{user} apa saja yang dikehendaki terlibat dalam sistem yang dirancang. Di samping itu, peserta wajib mencantumkan apa saja keterbatasan sistem yang dirancang, misalnya fitur apa saja yang tidak dapat atau tidak akan dilakukan atau tidak disediakan pada produk yang akan dikembangkan. 

    Pada bidang TE yang menghasilkan perangkat keras, dapat dijelaskan mengenai fitur fungsional dan karakteristik dari perangkat keras tersebut. Sebagai contoh, misalkan sebuah rancangan robot diharapkan dapat digunakan untuk membantu orang tua yang memiliki keterbatasan. Dalam memenuhi tugas tersebut, robot dapat dispesifikasikan memiliki kriteria tertentu, misalkan, stabil pada kondisi medan yang homogen maupun heterogen (lapangan rumput, aspal, con-block). Selain, itu robot dapat dispesifikasikan memiliki bentuk yang kokoh dan rapi, pengkabelan (\textit{wiring}) tidak berantakan, tertutup dan aman. Perlu disebutkan juga keterbatasan robot tersebut, misalnya hanya membahas robot beroda dan tidak bisa bergerak \textit{omni-directional}, daya yang tersedia maksimal adalah 40 Watt, dan sebagainya. 

    Pada bidang TE yang menghasilkan program simulasi, maka peserta \textit{capstone} dapat menspesifikasikan jenis program yang akan dibuat. Misalkan, sebuah program simulasi yang dibuat untuk menjalankan algoritme dalam rangka menjawab sebuah permasalahan (tentu saja permasalahan yang akan dipecahkan telah dijelaskan terlebih dahulu). Di samping juga batasan-batasan dari algoritme yang akan dirancang. Di samping itu, peserta \textit{capstone} secara umum harus menspesifikasikan skenario yang dipakai dalam proses perancangan perangkat lunak diatas. Pengujian bisa mespesifikasikan metrik/indikator yang akan diuji. Pada bidang STL yang tidak menghasilkan perangkat keras maupun program simulasi lengkap, maka peserta \textit{capstone} wajib menspesifikasikan analisis kebutuhan dan detil permasalahan yang akan dijawab, beserta batasan-batasan yang diprediksi mungkin timbul beserta alasan-alasan yang logis kenapa batasan-batasan tersebut dipilih/ditetapkan.

    
\section{Detail Rancangan}
\label{sec:Detail_Rancangan}
    
Tuliskan detail rancangan anda untuk menyelesaikan masalah \textit{engineering} yang kompleks di sini. Anda harus mendetailkan rancangan sampai level yang dapat diimplementasikan


% BAB 05 : Pengujian, Improvement, dan Pembahasan
\chapter{\uppercase{Pengujian dan Pembahasan}}
\label{chap:Pengujian_Improvement_dan_Pembahasan}
\section{Pengujian dan Pembahasan}
\label{sec:pengujian}


Pada bagian ini mahasiswa melakukan pengujian terhadap rancangan yang diajukan di \ref{chap:Detail_Implementasi}. Mahasiswa harus menuliskan hal-hal diantaranya: 
\begin{itemize}
    \item Tujuan pengujian,
    \item Hal atau bagian apa yang diuji,
    \item Skenario pengujian,
    \item Analisis dari hasil pengujian. Analisis harus dikaitkan dengan ilmu-ilmu dasar (matematika, fisika, dan \textit{engineering science}) kenapa suatu fenomena terjadi. Analisis juga merupakan cara untuk memastikan rancangan yang diajukan di \ref{chap:Detail_Implementasi} memenuhi spesifikasi atau kriteria sukses yang disebutkan,
    \item Hal-hal penting lainnya.
\end{itemize}

\subsection{Skenario Pengujian 1 dan Analisis}
\label{subsec:Skenario1}
Sesuaikan jumlah skenario sesuai kebutuhan pengujian anda
\subsection{Skenario Pengujian 2 dan Analisis}
\label{subsec:Skenario2}
Sesuaikan jumlah skenario sesuai kebutuhan pengujian anda

\section{\textit{Improvement}}
\label{sec:improvement}
Di dalam Bab ini, Anda harus secara rinci menjelaskan kenapa perlu tambahan dan/atau revisi rancangan secara logis dan kronologis sehingga tim penilai akan melihat bahwa tambahan ini adalah sesuatu yang perlu dan wajar untuk dilakukan. 




% BAB 06 : Impact dari the proposed engineering design ke aspek sosial dan lingkungan
\chapter{\uppercase{Analisis Mengenai Pengaruh Solusi \textit{Engineering Design} }}
\label{chap:Impact_Engineering}
Pada bab-bab sebelumnya, mahasiswa sudah melakukan identifikasi permasalahan \textit{engineering} yang kompleks, melakukan studi pustaka mengenai bagaimana masalah tersebut atau masalah serupa diselesaikan oleh orang lain, mengajukan solusi sendiri, dan melakukan detail rancangan dan pengujian.

Pada bab ini, mahasiswa melakukan analisis mengenai dampak dari solusi yang diajukan dalam konteks global, ekonomis, lingkungan, dan sosial. Anda dapat melakukan prediksi atau analisis apa yang akan terjadi kalau implementasi anda diadopsi secara luas.

% BAB 07 : Kesimpulan dan Saran
\chapter{\uppercase{Kesimpulan dan Saran}}
\label{chap:Kesimpulan_Saran}

\section{Kesimpulan}
\label{sec:kesimpulan}

    Kesimpulan pada dokumen C-501 ini berupa ringkasan dari proses desain produk, program, atau perangkat lunak yang telah dibuat, beserta hasil pengujian produk tersebut.

\section{Saran}
\label{sec:saran}

    Sebagai penutup, tuliskanlah saran-saran yang menurut Anda masih dapat dilakukan untuk memperbaiki kualitas/performa dari produk yang telah Anda buat.

% ------------------------------------------------------------ %
% References

\begin{thebibliography}{}

    \bibitem{edwards_sliding_1998}
Edwards,~C. and Spurgeon,~S.~K., \emph{Sliding Mode Control: Theory and Applications}, London, UK, 1998

\bibitem{b1} 
G. Eason, B. Noble, and I. N. Sneddon, ``On certain integrals of Lipschitz-Hankel type involving products of Bessel functions,'' Phil. Trans. Roy. Soc. London, vol. A247, pp. 529--551, April 1955.

\bibitem{b2} 
J. Clerk Maxwell, A Treatise on Electricity and Magnetism, 3rd ed., vol. 2. Oxford: Clarendon, 1892, pp.68--73.

\bibitem{b3} 
I. S. Jacobs and C. P. Bean, ``Fine particles, thin films and exchange anisotropy,'' in Magnetism, vol. III, G. T. Rado and H. Suhl, Eds. New York: Academic, 1963, pp. 271--350.

\bibitem{mics5524}
SGX Sensortech, ``MiCS-5524 Data Sheet,'' Rev. 6 (1084), Corcelles-Cormondr\`{e}che, Switzerland.

% TODO: lengkapi nama penulis (belum lengkap di daftar-pustaka.md [6])
\bibitem{mox_drift}
[Penulis perlu dilengkapi] \textit{et al.}, ``Correction model for metal oxide sensor drift caused by ambient temperature and humidity,'' \emph{Sensors}, vol. 22, no. 9, Art. no. 3301, 2022, doi: 10.3390/s22093301.


\end{thebibliography}

% Akhir Dokumen
\end{document}